\section{Extended model}
\label{sec:extended}

The extended model retains population growth, the integrated AI monopolist, the
nested production structure, the CES research aggregator, and the ideas production
function from Section~\ref{sec:toy}. It also retains
Assumption~\ref{ass:stable-human-essential-benchmark}. The model adds household
heterogeneity and richer accumulation of physical and computational capital. Regulated access to AI services
is treated as a policy counterfactual, not as a separate baseline industry.

\subsection{Households}

Each household type belongs to a population that grows at the common rate $n$.
Workers, researchers, and shareholders choose consumption and saving and own
specified shares of physical capital and firm equity. Workers receive production
wages, researchers receive research wages, and shareholders receive the integrated
developer's net profits.

This heterogeneity separates aggregate welfare from its distribution. Monopoly rents
can raise shareholder consumption while a high production elasticity $\sigma_{XL}$
between labor and AI services reduces labor income. The model will report aggregate
consumption, per-capita consumption, factor-income shares, and welfare by household type.

\subsection{Final-good firms}

Competitive final-good firms retain the production structure
\eqref{eq:final-production}--\eqref{eq:service-composite}. The substitution elasticity
$\sigma_{XL}$ determines how strongly firms replace production labor with AI services when
their relative price falls. Physical capital remains Cobb--Douglas with the CES
service composite. Final-good firms rent capital, hire production labor, and purchase
AI services at market prices.

\subsection{Integrated AI monopolist}

The integrated developer uses frontier capability $A_t$ and inference compute $U_t$
to supply quality-adjusted services $X_t$. It recognizes the demand generated by
final-good firms and chooses its price or quantity accordingly. Its operating rents
are the private return that links current market power to frontier research.

The baseline retains the normalization $X_t=A_tU_t$. A robustness exercise will allow
capability to affect inference efficiency through a separate mapping $B(A_t)$. This
will test whether the quantitative results depend on assigning the same state
variable to broad capability and services produced per unit of inference compute.

Inference and research compute are distinguished by use. The extended model can make
their resource cost depend on a common stock of computational capital and electricity,
without imposing an arbitrary fixed cap. This creates an endogenous opportunity cost
between current deployment and frontier research.

\subsection{Frontier research}

The integrated developer chooses human researchers and research compute,
\[
    \{H_t,M_t\},
\]
to maximize its market value. Capability evolves according to
\eqref{eq:ideas}--\eqref{eq:effective-research}. Human and automated-research
services are used simultaneously. Their equilibrium shares depend on the wage, the
common resource cost, and the substitution elasticity $\sigma_{HM}$. The private value of
an additional unit of capability reflects future cost reductions in supplying AI
services and the operating rents generated by those reductions.

\subsection{Government and social planner}

The government can regulate the access price for AI services, subsidize research
inputs, or pay an innovation prize. Net revenues are returned to households through
lump-sum transfers. The social planner values the current productive surplus from AI
services and the future knowledge benefits of capability. Its shadow value $Q_t$ can
therefore exceed the developer's private value $q_t$.

The planner's allocation provides the welfare benchmark. Comparing its conditions
with those of the monopolist identifies separately the static service-market markup
and the dynamic knowledge and appropriability wedge. The policy instruments are
specified in Section~\ref{sec:regulation}.

\subsection{Equilibrium}

Given a policy rule, an equilibrium consists of paths for population, consumption and
saving by household type, physical and computational capital, investment, production
and research labor, AI services, inference and research compute, and frontier
capability. Prices include the return to capital, wages, the price of AI services, the
common price of standardized compute when its supply is endogenous, and the private
shadow value of capability.

Households optimize given prices, profits, taxes, and transfers. Final-good firms
maximize profits. The integrated developer chooses current AI supply and frontier
research to maximize its market value subject to the ideas production function and
the policy regime. Labor and goods markets clear, and the laws of motion for
population, capital, compute capacity, and capability hold.

\subsection{Solution}

The solution traces a continuous change in the research-input composition. The
automated share is solved jointly with wages, investment, automated-research
services, and the private value of capability. Population growth anchors the
human-dominated benchmark. In the AI-dominated limit, capability growth depends on
the growth of automated-research services and on $\phi$ and $\eta$.
The solution must keep the two acceleration channels separate. It will report the
path of the research-feedback denominator as the automated contribution changes and
will distinguish that path from the production-side effect of $\sigma_{XL}>1$ on
the output--capital ratio and the return to capital.

For each policy regime, the model solves household and firm optimality conditions
together with market clearing and the state equations. The laissez-faire monopoly
and the regulated allocations are computed from the same initial conditions and
technological parameters. The social-planner allocation then decomposes the welfare
gap into restricted current deployment, imperfect appropriation of innovation
benefits, and distributional effects.

The quantitative solution will report both the production-labor share,
$w_tL_t/Y_t$, and the aggregate labor-income share, $w_tN_t/Y_t$. The first
isolates substitution between production labor and AI services. The second also
incorporates the endogenous allocation of labor to frontier research. This
distinction is necessary for taking Proposition
\ref{prop:wage-share-thresholds} from a conditional production result to a
general-equilibrium distributional result.

\subsubsection{Quantitative objects}

The solution requires:
\begin{enumerate}[label=(\roman*)]
    \item population growth $n$ and the allocation of labor between production and
    research;
    \item the substitution elasticity $\sigma_{XL}$ between production labor and AI
    services, together with the CES weights $(\omega_L,\omega_X)$;
    \item the common unit resource cost $\xi$ of inference and research compute and,
    in the richer version, the law of motion for computational capital;
    \item the research-substitution elasticity $\sigma_{HM}$ and the CES weights
    $(\omega_H,\omega_M)$;
    \item the elasticities $\phi$ and $\eta$ in the ideas production function;
    \item a capability-dependent mapping $\lambda(A)$ for automated research in a
    robustness extension;
    \item the mapping $B(A)$ from capability to inference efficiency in the
    robustness exercise; and
    \item the degree to which operating rents and innovation prizes let the
    developer appropriate the social value of capability.
\end{enumerate}
These objects determine human-dominated growth, the speed of research automation,
whether recursive feedback is compatible with constant growth, whether production
admits a finite-growth limit, the output cost of monopoly pricing, and the gap
between private and social research.
